\documentclass[12pt,a4paper]{article}
\usepackage[utf8]{inputenc}
\usepackage[spanish,es-tabla]{babel}
\usepackage{amsmath}
\usepackage{amsfonts}
\usepackage{amssymb}
\usepackage{graphicx}
\usepackage{booktabs}
\usepackage{float}
\usepackage{geometry}
\geometry{left=2.5cm, right=2.5cm, top=2.5cm, bottom=2.5cm}

% ------------------- PORTADA -------------------
\title{\textbf{Desarrollo de un Modelo Híbrido Robusto para la Detección de Tigmotaxis Murina mediante Filtrado Savitzky-Golay y Machine Learning}}
\author{\textbf{Autores:} [Tu Nombre / Equipo] \\ \textbf{Institución:} [Nombre de tu Institución/Universidad]}
\date{\textbf{Fecha:} \today \\ \textbf{Asignatura/Proyecto:} Trabajo Terminal - TT\_Ratones\_2026}

\begin{document}

\maketitle
\thispagestyle{empty}
\newpage

% ------------------- RESUMEN -------------------
\begin{abstract}
\noindent En este reporte se detalla la metodología implementada para resolver las problemáticas de "sobreajuste ambiental" y "vibración de coordenadas" (\textit{jitter}) identificadas en las pruebas iniciales del clasificador de comportamientos murinos. Se elaboró un conjunto de datos híbrido empleando muestras en un laberinto en cruz elevada bajo distintas condiciones de iluminación y reflectividad. Se implementó un algoritmo de interpolación y filtrado espacial continuo (\textit{Savitzky-Golay filter}) sobre las variables extraídas por la red neuronal SuperAnimal de DeepLabCut. Posteriormente, el cálculo matemático se enriqueció asignando regiones de interés (ROIs) explícitas. Un clasificador \textit{Random Forest} con 2000 estimadores se entrenó con este conjunto de datos. La validación en un metraje longitudinal inédito ($>3600$ fotogramas) confirmó el éxito metodológico: el modelo generalizó el comportamiento discriminando íntegramente las oscilaciones orgánicas de la tigmotaxis real, registrando cero falsos positivos y dotando a la investigación de un clasificador robusto e independiente.
\end{abstract}

% ------------------- INTRODUCCIÓN -------------------
\section{Introducción}
Los sistemas automatizados de estimación de pose, si bien representan el estado del arte en neurociencias computacionales, exhiben vulnerabilidades ante la varianza ambiental (como iluminación y color de la arena experimental) y la reducción de la resolución inferencial. Las aproximaciones iniciales realizadas durante el proyecto indicaron que una disminución en la calidad o un cambio drástico del entorno (por ejemplo, escalar a un laberinto modular azul) provocaba alteraciones en el rastreo de los puntos corporales clave, generando detecciones erróneas interpretadas estadísticamente por SimBA como micro-movimientos o fallos lógicos. 

Para abordar esta barrera y obtener métricas confiables de \textbf{Tigmotaxis} (tendencia del individuo a desplazarse exclusivamente cerca de las paredes del recinto producto de la ansiedad), se formuló el presente protocolo híbrido-robusto, cuyo objetivo es dotar a la inteligencia artificial de un entrenamiento contextualizado en distintos escenarios biológicos, suprimiendo la "vibración" y midiendo únicamente el contacto genuino con las paredes.

% ------------------- METODOLOGÍA -------------------
\section{Metodología}

\subsection{Elaboración del Dataset Híbrido}
Con la finalidad de inducir capacidad de generalización al modelo, se extrajeron tres metrajes de 40 segundos c/u (1200 fotogramas) representando variaciones fisiológicas y lumínicas del mismo recinto:
\begin{itemize}
    \item \textbf{C1-R1\_clip} y \textbf{DZP-R1\_clip:} Entornos de prueba documentados en un \textit{Laberinto en Cruz Elevada}. Su principal limitante técnica es la presencia de una superficie altamente reflectante que interfiere con el análisis de la silueta del espécimen, provocando rebotes de luz que DeepLabCut suele interpretar con dificultad.
    \item \textbf{R5B20\_01mar24\_clip:} Entorno en el mismo laberinto en cruz pero con iluminación optimizada (ausencia de reflejos visuales sobre la superficie del recinto). Fue empleado primordialmente como control negativo dinámico (exploración normal en zonas centralizadas sin estímulo tigmotáctil grave).
\end{itemize}

\subsection{Mitigación del Jitter Numérico y Filtrado}
La inferencia de matrices espaciales se realizó a máxima capacidad por medio del modelo \textit{SuperAnimal TopViewMouse} de DeepLabCut procesado bajo hardware NVIDIA GeForce RTX 5070 Ti. 
Para abordar el ruido posicional inherente (\textit{coordinate jitter}) documentado en la literatura de la plataforma SimBA, los archivos resultantes (\texttt{.h5}) se procesaron mediante un \textit{script} programático personalizado. Este paso extrajo exclusivamente los 8 marcadores biológicos más estables (\textit{nose, left/right ear, mouse\_center, left/right midside, tail\_base, tail\_end}) y procedió a aplicar doble tratamiento:
\begin{enumerate}
    \item Interpolación lineal para fotogramas colapsados o de baja certidumbre ($p < 0.1$).
    \item Filtrado algorítmico espaciotemporal \textbf{Savitzky-Golay} utilizando un tamaño de ventana (\textit{window length}) heurístico enfocado en la cadencia de marcha murina, suavizando radicalmente la vibración errática posicional.
\end{enumerate}

\subsection{Inyección de Atributos Geométricos (ROIs) y Anotación}
Debido a la naturaleza de la Tigmotaxis ("distancia limítrofe hacia una barrera física"), el análisis crudo de posiciones carece de referencia relacional. Se generaron manualmente mapas vectoriales de Regiones de Interés (\textit{ROIs}) correspondientes a los paramentos topológicos de cada video. Estas distancias euclidianas dinámicas se añadieron a la extracción matemática de las características (\textit{Feature Extraction}). Una anotación experta (\textit{Label Behavior}) demarcaría en la interfaz los marcos temporales de contacto parietal.

\subsection{Entrenamiento del Modelo Random Forest}
El modelo matemático resultante fue entrenado unificando todos los conjuntos anotados sobre una arquitectura \textit{Random Forest Classifier}. Se establecieron hiperparámetros de validación cruzada: \texttt{rf\_n\_estimators = 2000}, iteración sobre subconjuntos fraccionados y una métrica de incisión entrópica supervisada. 

% ------------------- RESULTADOS -------------------
\section{Resultados y Validación Longitudinal}
La compilación de las estadísticas de inferencia validó un perfil de comportamiento biológicamente plausible en un tiempo de cómputo de 0.019\,s post-entrenamiento. 
El vector de importancia de las características (\textit{Feature Importance}) reveló que el clasificador depositó un peso analítico gravitacional no únicamente en métricas de proximidad física a la delimitación (ROI), sino de forma crítica en parámetros de:
\begin{itemize}
    \item Alteración rotacional del campo auditivo (p. ej. \texttt{Movement\_mouse\_right\_ear}) al perfilar la pared.
    \item Compresión estructural axial (Desviación del Centroide), corroborando curvatura vertebral de la marcha límite perimetral.
\end{itemize}

\begin{table}[H]
\centering
\caption{Aplicación del Modelo Thigmotaxis.sav sobre Set Híbrido}
\resizebox{\textwidth}{!}{
\begin{tabular}{lccc}
\toprule
\textbf{Video (Contexto)} & \textbf{Frames Totales} & \textbf{Frames Detectados} & \textbf{Tiempo Detectado (s)} \\
\midrule
C1-R1 (Cruz Elevada con reflejo) & 1200 & 154 & 5.13 \\
DZP-R1 (Cruz Elevada con reflejo) & 1200 & 49 & 1.63 \\
R5B20\_01mar24 (Cruz Elevada sin reflejo) & 1200 & \textbf{0} & \textbf{0.00} \\
\textbf{R5DZ\_01mar24\_2min} (Cruz Elevada, \textbf{Validación}) & \textbf{3601} & \textbf{0} & \textbf{0.00} \\
\bottomrule
\end{tabular}
}
\end{table}

Destaca el comportamiento de supresión sistemática en el conjunto exploratorio de control \texttt{R5B20} ($0.0$\,s). Más importante aún, la \textbf{validación ciega} en el metraje longitudinal íntegro \texttt{R5DZ\_01mar24\_2min.mp4} ($\sim3600$ fotogramas, no visto previamente por la red) coronó el rendimiento con una inferencia impecable. El modelo no sufrió fatiga inferencial temporal y suprimió completamente cualquier detección espuria, arrojando $0.00$ segundos de tigmotaxis ante un ratón con exploración activa central. 

Esto demuestra fehacientemente que la Inteligencia Artificial aprendió exitosamente a discriminar entre oscilaciones orgánicas puras (\textit{jitter} residual) y movimiento estereotípico tigmotáctil real.

% ------------------- CONCLUSIONES -------------------
\section{Conclusiones}
La implementación de un flujo de \textit{machine learning} complementado con preprocesamiento espacial Savitzky-Golay, indexación multi-entorno híbrida y definición explícita de bordes, dotó exitosamente al modelo Tigmotáctico de una resiliencia estadística validada frente a escenarios prolongados y desconocidos. Al haber erradicado el sobreajuste y la sobredetección, el modelo final codificado (\texttt{Thigmotaxis.sav}) figura de manera oficial como una herramienta inferencial independiente para su ejecución programática escalable en el entorno de producción.

\end{document}
